% !TeX root = main.tex

\setlength{\abovedisplayskip}{7pt}
\setlength{\belowdisplayskip}{7pt}
\setlength{\abovedisplayshortskip}{4pt}
\setlength{\belowdisplayshortskip}{4pt}

% Места для текста и рисунков заполняются после моделирования.
\newcommand{\reporttodo}[1]{%
  \par\smallskip\noindent\textit{[Заполнить: #1]}\par\smallskip
}

% Рисунок и подпись располагаются рядом с соответствующим текстом.
\makeatletter
\newenvironment{reportfigure}{%
  \par\addvspace{8pt}\noindent
  \begin{minipage}{\textwidth}
    \def\@captype{figure}%
    \centering
}{%
  \end{minipage}\par\vspace{4pt}
}
\makeatother

% Аргументы: файл в images, подпись, метка. После добавления файла
% рамка автоматически заменяется изображением.
\newcommand{\imageplaceholder}[3]{%
  \begin{reportfigure}
    \IfFileExists{images/#1}{%
      \includegraphics[width=\textwidth,height=0.7\textheight,keepaspectratio]{#1}%
    }{%
      \fbox{\parbox[c][3cm][c]{0.9\textwidth}{%
        \centering Место для рисунка\\[0.4em]
        \texttt{images/\detokenize{#1}}%
      }}%
    }
    \caption{#2}
    \label{#3}
  \end{reportfigure}
}

\tableofcontents
\newpage

\section{Цель работы}

Получить дискретные алгоритмы ПИД-регулирования и исследовать влияние
метода аппроксимации, периода дискретизации и неточности модели объекта
на переходные процессы при наличии шума измерения.
Построить модель нерегулярной качки и синтезировать фильтры Калмана.

\section{Непрерывный ПИД-регулятор}

\subsection{Модель объекта и исходные данные}

Для варианта 3 заданы постоянные времени
\begin{equation}
  T_1=0{,}6\,\text{с},\qquad T_2=0{,}9\,\text{с}
  \label{eq:plant-data}
\end{equation}
Объект состоит из двух последовательно соединённых апериодических звеньев
с единичными коэффициентами передачи:
\begin{equation}
  W_0(s)=\frac{1}{(T_1s+1)(T_2s+1)}
    =\frac{1}{0{,}54s^2+1{,}5s+1}
  \label{eq:plant-transfer}
\end{equation}
При обозначении задания через $r$, истинного выхода через $y$
и шума измерения через $v$ получаем
\begin{equation}
  y_{\text{и}}(k)=y(k)+v(k),
  \qquad e(k)=r(k)-y_{\text{и}}(k)
  \label{eq:measurement-error}
\end{equation}

\subsection{Компенсация полюсов объекта}

Непрерывный регулятор записывается в виде
\begin{equation}
  W_{\text{ПИД}}(s)=q_0\left(k_p+\frac{k_i}{s}+k_ds\right)
  \label{eq:continuous-controller}
\end{equation}
Здесь $q_0$ является общим коэффициентом усиления регулятора.
Для компенсации полюсов объекта выбираем
\begin{equation}
  k_p=T_1+T_2=1{,}5,
  \qquad k_i=1,
  \qquad k_d=T_1T_2=0{,}54
  \label{eq:pid-gains}
\end{equation}
Передаточная функция регулятора принимает вид
\begin{equation}
  W_{\text{ПИД}}(s)
  =q_0\frac{0{,}54s^2+1{,}5s+1}{s}
  =q_0\frac{(0{,}6s+1)(0{,}9s+1)}{s}
  \label{eq:continuous-controller-numeric}
\end{equation}
Для номинального объекта передаточные функции разомкнутой
и замкнутой систем равны
\begin{equation}
  W_{\text{ПИД}}(s)W_0(s)=\frac{q_0}{s},
  \qquad W_{\text{зам}}(s)=\frac{q_0}{s+q_0}
  \label{eq:continuous-closed-loop}
\end{equation}

\section{Дискретные алгоритмы ПИД-регулирования}

\subsection{Рекуррентные уравнения}

В реализуемых алгоритмах производная ошибки аппроксимируется обратной
разностью, а способы интегрирования различаются. Обозначим через $I(k)$
интегральную составляющую без общего усиления регулятора:
\begin{equation}
  u(k)=q_0\left[k_pe(k)+I(k)
    +\frac{k_d}{T}\bigl(e(k)-e(k-1)\bigr)\right]
  \label{eq:pid-components}
\end{equation}
При прямом методе Эйлера используется предыдущее значение ошибки:
\begin{equation}
  I(k)=I(k-1)+k_iT e(k-1)
  \label{eq:forward-integral}
\end{equation}
Для обратного метода Эйлера имеем
\begin{equation}
  I(k)=I(k-1)+k_iT e(k)
  \label{eq:backward-integral}
\end{equation}
При методе трапеций, соответствующем интегрированию по Тастину,
получаем
\begin{equation}
  I(k)=I(k-1)+\frac{k_iT}{2}\bigl(e(k)+e(k-1)\bigr)
  \label{eq:tustin-integral}
\end{equation}
После вычитания соседних значений управления все три алгоритма
принимают форму
\begin{equation}
  u(k)=u(k-1)+q_0\left[c_0e(k)+c_1e(k-1)+c_2e(k-2)\right]
  \label{eq:incremental-pid}
\end{equation}
Коэффициенты $c_0,c_1,c_2$ определяют веса отсчётов ошибки;
общий коэффициент $q_0$ вынесен отдельно.

\begin{table}[!htbp]
  \centering
  \caption{Коэффициенты рекуррентного ПИД-регулятора}
  \label{tab:pid-general}
  \begin{tabular}{c c c c}
    \toprule
    Метод & $c_0$ & $c_1$ & $c_2$\\
    \midrule
    Прямой Эйлер & $k_p+k_d/T$ & $-k_p+k_iT-2k_d/T$ & $k_d/T$\\
    Обратный Эйлер & $k_p+k_iT+k_d/T$ & $-k_p-2k_d/T$ & $k_d/T$\\
    Трапеции & $k_p+k_iT/2+k_d/T$ & $-k_p+k_iT/2-2k_d/T$ & $k_d/T$\\
    \bottomrule
  \end{tabular}
\end{table}

\subsection{Численные коэффициенты}

Периоды дискретизации для сравнения равны
\begin{equation}
  T=\frac{T_1}{2}=0{,}3\,\text{с},
  \qquad T=\frac{T_1}{4}=0{,}15\,\text{с}
  \label{eq:sample-periods}
\end{equation}
Подстановка коэффициентов~\eqref{eq:pid-gains} даёт значения,
приведённые в таблице~\ref{tab:pid-numeric}.

\begin{table}[!htbp]
  \centering
  \caption{Численные веса отсчётов ошибки}
  \label{tab:pid-numeric}
  \begin{tabular}{c c r r r}
    \toprule
    $T$, с & Метод & $c_0$ & $c_1$ & $c_2$\\
    \midrule
    $0{,}3$ & Прямой Эйлер & $3{,}3$ & $-4{,}8$ & $1{,}8$\\
    $0{,}3$ & Обратный Эйлер & $3{,}6$ & $-5{,}1$ & $1{,}8$\\
    $0{,}3$ & Трапеции & $3{,}45$ & $-4{,}95$ & $1{,}8$\\
    \midrule
    $0{,}15$ & Прямой Эйлер & $5{,}1$ & $-8{,}55$ & $3{,}6$\\
    $0{,}15$ & Обратный Эйлер & $5{,}25$ & $-8{,}7$ & $3{,}6$\\
    $0{,}15$ & Трапеции & $5{,}175$ & $-8{,}625$ & $3{,}6$\\
    \bottomrule
  \end{tabular}
\end{table}

\section{Исследование замкнутой системы}

\subsection{Подбор усиления и сравнение методов}

В первом опыте используется период $T=0{,}3\,\text{с}$,
а дисперсия аддитивного шума измерения задана условием
\begin{equation}
  \operatorname{Var}v(k)=0{,}01
  \label{eq:pid-measurement-variance}
\end{equation}
Коэффициент $q_0$ подбирается экспериментально для устойчивого
слабоколебательного процесса. В последующих сравнениях для каждого
способа сохраняется найденное для него усиление.

\reporttodo{указать подобранные значения $q_0$, параметры ступенчатого
задания, начальные состояния, интервал моделирования и условия генерации шума}

\imageplaceholder{task2_pid_methods.png}
  {Процессы в замкнутой системе при трёх способах аппроксимации и $T=0{,}3$ с}
  {fig:pid-methods}

\reporttodo{сравнить переходные процессы, их колебательность,
длительность и чувствительность к шуму измерения}

\subsection{Изменение периода дискретизации}

Для второго опыта принимается $T=0{,}15\,\text{с}$, причём коэффициенты
непрерывного ПИД-регулятора и подобранное усиление сохраняются,
а веса дискретного алгоритма берутся из таблицы~\ref{tab:pid-numeric}.

\imageplaceholder{task3_sample_period.png}
  {Процессы при уменьшении периода дискретизации до $T=0{,}15$ с}
  {fig:sample-period}

\reporttodo{сопоставить результаты при двух периодах дискретизации
для каждого метода при одинаковом усилении и шуме измерения}

\subsection{Неточность компенсации полюсов}

Постоянную времени объекта изменяем относительно номинального значения:
\begin{equation}
  T_2^-=0{,}8T_2=0{,}72\,\text{с},
  \qquad T_2^+=1{,}2T_2=1{,}08\,\text{с}
  \label{eq:uncertain-time-constants}
\end{equation}
Передаточные функции изменённого объекта равны
\begin{equation}
  W_0^-(s)=\frac{1}{0{,}432s^2+1{,}32s+1},
  \qquad
  W_0^+(s)=\frac{1}{0{,}648s^2+1{,}68s+1}
  \label{eq:uncertain-plants}
\end{equation}
Коэффициенты регулятора сохраняются рассчитанными для номинального
объекта; период и усиление принимаются из первого опыта.

\imageplaceholder{task4_pole_compensation.png}
  {Сравнение процессов при номинальном и изменённых значениях $T_2$}
  {fig:pole-compensation}

\reporttodo{оценить влияние увеличения и уменьшения $T_2$ на качество
процессов при неизменных коэффициентах регулятора}

\section{Модель нерегулярной качки}

По данным варианта 3 параметры модели и шумов равны
\begin{equation}
  \omega_0=0{,}9,\qquad \alpha=0{,}05,\qquad G_0=2,
  \qquad Q=1{,}5,\qquad R=0{,}8
  \label{eq:pitching-data}
\end{equation}
Формирующий фильтр задаётся передаточной функцией
\begin{equation}
  \Phi(s)=\frac{\omega_0^2}{s^2+2\alpha s+\omega_0^2}
    =\frac{0{,}81}{s^2+0{,}1s+0{,}81}
  \label{eq:pitching-transfer}
\end{equation}
Параметр $G_0$ определяет уровень спектральной плотности белого шума
на входе непрерывного формирующего фильтра.
Для дискретной модели дисперсия входного шума задаётся отдельно
величиной $Q$ из таблицы исходных данных.

При выборе состояний $\bm x=(y,\dot y)^{\mathsf T}$ непрерывные матрицы равны
\begin{equation}
  A_{\text{к}}=\begin{pmatrix}0&1\\-0{,}81&-0{,}1\end{pmatrix},
  \qquad
  B_{\text{к}}=\begin{pmatrix}0\\0{,}81\end{pmatrix},
  \qquad C_{\text{к}}=\begin{pmatrix}1&0\end{pmatrix}
  \label{eq:pitching-continuous-matrices}
\end{equation}
Принимаем период дискретизации $h=0{,}1\,\text{с}$ и используем
точную дискретизацию с удержанием входа:
\begin{equation}
  A_h=e^{A_{\text{к}}h},
  \qquad
  B_h=A_{\text{к}}^{-1}(A_h-I_2)B_{\text{к}}
  \label{eq:pitching-discretization}
\end{equation}
Численные матрицы равны
\begin{equation}
  A_h\approx\begin{pmatrix}
    0{,}99596619&0{,}09936739\\
    -0{,}08048759&0{,}98602945
  \end{pmatrix},
  \qquad
  B_h\approx\begin{pmatrix}0{,}00403381\\0{,}08048759\end{pmatrix}
  \label{eq:pitching-discrete-matrices}
\end{equation}
Модель с шумами записывается в виде
\begin{equation}
  \bm x(k+1)=A_h\bm x(k)+B_hw(k),
  \qquad y_{\text{и}}(k)=C_{\text{к}}\bm x(k)+v(k)
  \label{eq:pitching-discrete-model}
\end{equation}
Шумы $w$ и $v$ имеют нулевые средние, не коррелированы друг с другом
и удовлетворяют условиям
\begin{equation}
  \operatorname{Var}w(k)=Q=1{,}5,
  \qquad \operatorname{Var}v(k)=R=0{,}8
  \label{eq:pitching-noise-variance}
\end{equation}

\reporttodo{указать начальное состояние, интервал моделирования
и условия генерации шумов}

\imageplaceholder{task5_pitching.png}
  {Истинный процесс нерегулярной качки и зашумлённое измерение}
  {fig:pitching-model}

\reporttodo{описать полученный случайный процесс и влияние шума измерения}

\section{Фильтр Калмана для модели качки}

Ковариация возмущения в уравнении состояния равна
\begin{equation}
  Q_x=B_hQB_h^{\mathsf T}
  \label{eq:pitching-state-covariance}
\end{equation}
Обозначим через $P^-$ стационарную ковариацию ошибки прогноза.
Она находится из дискретного уравнения Риккати
\begin{equation}
  \begin{aligned}
    P^-={}&A_hP^-A_h^{\mathsf T}+Q_x\\
      &-A_hP^-C_{\text{к}}^{\mathsf T}
        \bigl(C_{\text{к}}P^-C_{\text{к}}^{\mathsf T}+R\bigr)^{-1}
        C_{\text{к}}P^-A_h^{\mathsf T}
  \end{aligned}
  \label{eq:pitching-riccati}
\end{equation}
Коэффициент коррекции текущей оценки определяется выражением
\begin{equation}
  K_{\text{к}}=P^-C_{\text{к}}^{\mathsf T}
    \bigl(C_{\text{к}}P^-C_{\text{к}}^{\mathsf T}+R\bigr)^{-1}
  \label{eq:pitching-kalman-gain}
\end{equation}
Решение уравнения~\eqref{eq:pitching-riccati} и коэффициент фильтра равны
\begin{equation}
  P^-\approx\begin{pmatrix}
    0{,}08111358&0{,}03913941\\
    0{,}03913941&0{,}10320731
  \end{pmatrix},
  \qquad
  K_{\text{к}}\approx\begin{pmatrix}0{,}09205803\\0{,}04442039\end{pmatrix}
  \label{eq:pitching-kalman-numeric}
\end{equation}
Алгоритм включает коррекцию по текущему измерению и прогноз
на следующий отсчёт:
\begin{equation}
  \begin{aligned}
    \hat{\bm x}_{k|k}
      &=\hat{\bm x}_{k|k-1}
        +K_{\text{к}}\bigl(y_{\text{и}}(k)-C_{\text{к}}\hat{\bm x}_{k|k-1}\bigr),\\
    \hat{\bm x}_{k+1|k}&=A_h\hat{\bm x}_{k|k},\\
    \hat y(k)&=C_{\text{к}}\hat{\bm x}_{k|k}
  \end{aligned}
  \label{eq:pitching-kalman-algorithm}
\end{equation}

\reporttodo{указать начальную оценку состояния и условия опыта}

\imageplaceholder{task7_pitching_kalman.png}
  {Процесс качки, его измерение и оценка фильтра Калмана}
  {fig:pitching-kalman}

\reporttodo{сравнить истинный сигнал, зашумлённое измерение и оценку;
оценить ошибку восстановления по результатам моделирования}

\section{Фильтр Калмана в контуре ПИД-регулирования}

Для нечётного варианта 3 используем ПИД-алгоритм прямого Эйлера
при периоде $T=0{,}3\,\text{с}$ из первого опыта.
Для расчёта фильтра выберем состояния объекта
$\bm x=(y,\dot y)^{\mathsf T}$ и запишем непрерывные матрицы:
\begin{equation}
  A_0=\begin{pmatrix}0&1\\-1/0{,}54&-1{,}5/0{,}54\end{pmatrix},
  \qquad B_0=\begin{pmatrix}0\\1/0{,}54\end{pmatrix},
  \qquad C=\begin{pmatrix}1&0\end{pmatrix}
  \label{eq:temperature-continuous-matrices}
\end{equation}
При удержании управления между отсчётами точная дискретная модель
объекта имеет матрицы
\begin{equation}
  \begin{aligned}
    A_{\text{д}}&=e^{A_0T}
      \approx\begin{pmatrix}
        0{,}93653261&0{,}19800117\\
        -0{,}36666884&0{,}38652936
      \end{pmatrix},\\
    B_{\text{д}}&=A_0^{-1}(A_{\text{д}}-I_2)B_0
      \approx\begin{pmatrix}0{,}06346739\\0{,}36666884\end{pmatrix}
  \end{aligned}
  \label{eq:temperature-discrete-matrices}
\end{equation}
Фильтр использует измеренный выход и известное фактически приложенное
управление. Расчётная модель имеет вид
\begin{equation}
  \bm x(k+1)=A_{\text{д}}\bm x(k)+B_{\text{д}}u(k)+\bm w_x(k),
  \qquad y_{\text{и}}(k)=C\bm x(k)+v(k)
  \label{eq:temperature-kalman-model}
\end{equation}
Дисперсия шума измерения для варианта 3 равна
\begin{equation}
  R_{\text{ПИД}}=0{,}03
  \label{eq:temperature-measurement-variance}
\end{equation}
При заданной ковариации $Q_{\text{ПИД}}$ возмущения состояния
ковариация прогноза и коэффициент коррекции находятся из уравнений
\begin{equation}
  \begin{aligned}
    P^-_{\text{ПИД}}={}&A_{\text{д}}P^-_{\text{ПИД}}A_{\text{д}}^{\mathsf T}
      +Q_{\text{ПИД}}\\
      &-A_{\text{д}}P^-_{\text{ПИД}}C^{\mathsf T}
        \bigl(CP^-_{\text{ПИД}}C^{\mathsf T}+R_{\text{ПИД}}\bigr)^{-1}
        CP^-_{\text{ПИД}}A_{\text{д}}^{\mathsf T},\\
    K_{\text{ПИД}}={}&P^-_{\text{ПИД}}C^{\mathsf T}
      \bigl(CP^-_{\text{ПИД}}C^{\mathsf T}+R_{\text{ПИД}}\bigr)^{-1}
  \end{aligned}
  \label{eq:temperature-riccati}
\end{equation}
Уравнения фильтра:
\begin{equation}
  \begin{aligned}
    \hat{\bm x}_{k|k}
      &=\hat{\bm x}_{k|k-1}
        +K_{\text{ПИД}}\bigl(y_{\text{и}}(k)-C\hat{\bm x}_{k|k-1}\bigr),\\
    \hat{\bm x}_{k+1|k}&=A_{\text{д}}\hat{\bm x}_{k|k}+B_{\text{д}}u(k),\\
    \hat y(k)&=C\hat{\bm x}_{k|k}
  \end{aligned}
  \label{eq:temperature-kalman-algorithm}
\end{equation}
В регулятор поступает ошибка $e(k)=r(k)-\hat y(k)$;
коэффициенты ПИД-регулятора сохраняются из опыта с прямым методом Эйлера.

Ковариация $Q_{\text{ПИД}}$ оставляется параметром для этапа моделирования.

\reporttodo{выбрать ковариацию $Q_{\text{ПИД}}$ и рассчитать численный
коэффициент фильтра; указать начальную оценку и условия опыта;
для сравнения с фильтром и без него использовать одинаковую дисперсию
шума измерения $R_{\text{ПИД}}=0{,}03$ в обоих опытах}

\imageplaceholder{task8_pid_kalman.png}
  {Переходные процессы в ПИД-системе с фильтром Калмана и без него}
  {fig:pid-kalman}

\reporttodo{сравнить качество переходного процесса и влияние шума
измерения при использовании фильтра и без фильтра}

\section{Выводы}

\reporttodo{по результатам моделирования сопоставить методы
аппроксимации, влияние периода дискретизации и ошибки компенсации
полюсов; оценить работу фильтров Калмана}
